\section{Certified companion identities}\label{sec:cert}

\subsection{The second companion}

\begin{theorem}\label{thm:k2}
For every $m\ge2$, identity \eqref{eq:wron} holds for $k=2$ and $z=0,\frac12$. Consequently, the three pairs $(B(2,z),C(2,z))$ with $z=-\frac12,0,\frac12$ lie in $\M_{2m+1}$.
\end{theorem}

Let $v\in\Q$ and $u=q^v$, and set
\[
E_i=e_{2i}-e_{2i-1},\qquad G_{m,2}(v;q)=F_{2m+1,B_v+E_{m-1}+E_m,0}(q).
\]

\begin{lemma}\label{lem:k2red}
For every $m\ge2$,
\begin{multline}\label{eq:k2red}
G_{m,2}(v;q)=q^{-2v-2}(-q^{1/2-v};q)_\infty\Bigl(\HL_m(q^2;u\one;q)\\
-(1+q+q^2)\HL_m(q;u\one;q)+q\HL_m(1;u\one;q)\Bigr).
\end{multline}
\end{lemma}

\begin{proof}
Only the last five coordinates are used. Introduce indeterminates $a,b,c,d_0,e$ corresponding to $q^{n_{r-4}},\dots,q^{n_r}$, respectively, and put $M=abcd_0e$, $K=1+w/u$ and $R=Ke-1$. For a Laurent polynomial $P$, let $\sigma_bP$ be obtained by replacing $b$ by $qb$, and define $\sigma_c,\sigma_{d_0},\sigma_e$ similarly. Consider the four operators
\begin{align*}
\mathcal A_2P&=ac(1-b)P-u^2b^2\sigma_bP,\\
\mathcal A_3P&=bd_0(1-c)P-u^{-2}c^2\sigma_cP,\\
\mathcal A_4P&=ce(1-d_0)P-u^2d_0^2\sigma_{d_0}P,\\
\mathcal A_5P&=d_0(1-e)P-(wu)^{-1}e\,\sigma_eP,
\end{align*}
and the four Laurent polynomials
\begin{align*}
P_2&=d_0e\bigl(Ke^2-(K(1+w^4)+w^2)e+w^2\bigr)+\frac{w^3}ue^2R-w^3ud_0^2eR,\\
P_3&=w^3ua\,d_0eR,\\
P_4&=a\left(\frac{w^3}u\,beR-\frac{w^4}{u^2}\,ce-w^4Ke+w^3uKd_0(1-e)\right),\\
P_5&=ace\left(Ke-w^2(w^2+wu+1)-\frac{w^3}u\,c\right)+u^2b^2e\bigl(w^3uKd_0+w^2-Ke\bigr)+w^5u^3Ka\,d_0^2.
\end{align*}
Direct expansion gives the Laurent-polynomial identity
\begin{multline}\label{eq:k2cert}
\sum_{i=2}^5\mathcal A_iP_i=w^4u^2b^2d_0^2e\\
-M\Bigl((1+w/u)(1+w^3/u)e^2-(1+w^2+w^4)(1+w/u)e+w^2\Bigr).
\end{multline}
The four displayed polynomials constitute the complete certificate.

For a Laurent polynomial $P$ in these five variables, let
\[
\Lambda(P)=\sum_{n\in\Z_{\ge0}^r}\frac{q^{n^\T T_rn/2+B_v^\T n}}{\prod_{i=1}^r(q;q)_{n_i}}\,P(q^{n_{r-4}},\dots,q^{n_r}).
\]
The reindexing used in \eqref{eq:contig} gives $\Lambda(M^{-1}\mathcal A_iP)=0$ for $i=2,3,4,5$. For example, the shift in $b$ has factor $q^{1+2v}$ before conjugation by $M^{-1}$; the latter supplies $q^{-1}$, leaving the factor $u^2$ displayed in $\mathcal A_2$. At the last coordinate the factor is $q^{1/2-v-1}=(wu)^{-1}$. The factors $1-q^{n_i}$ remove the lower boundary terms, and positive definiteness justifies all reindexings, including those with Laurent-polynomial weights.

Apply $\Lambda(M^{-1}\,\cdot\,)$ to \eqref{eq:k2cert}. Writing $F_j=F_{r,B_v+je_r,0}$ gives
\[
q^{2v+2}G_{m,2}(v;q)=(1+q^{1/2-v})(1+q^{3/2-v})F_2-(1+q+q^2)(1+q^{1/2-v})F_1+qF_0.
\]
\Cref{lem:boundary} expresses $F_j$ as $(-q^{j+1/2-v};q)_\infty\HL_m(q^j;u\one;q)$. Extracting $(-q^{1/2-v};q)_\infty$ proves the assertion.
\end{proof}

\begin{proof}[Proof of \cref{thm:k2}]
\Cref{lem:bdexp,lem:k2red} give
\begin{equation}\label{eq:k2final}
G_{m,2}(v;q)=q^{-2v}(-q^{1/2-v};q)_\infty\,J_2(q^v\one).
\end{equation}
At $v=0$ this is $S^{(0)}_{m,2}=(-w;q)_\infty J_2(\one)$, and \eqref{eq:conf0} gives the claimed Wronskian formula. At $v=\frac12$ the left-hand side is $S^{(-1/2)}_{m,2}$. Using \eqref{eq:halfshift} and $(-1;q)_\infty=2(-q;q)_\infty$ in \eqref{eq:k2final} yields
\[
S^{(1/2)}_{m,2}=q(-q;q)_\infty J_2(w\one).
\]
The factor $q^{-1}$ in \eqref{eq:confhalf} with $k=2$ cancels $q$, which proves the positive half-shift formula. Modularity follows from \cref{prop:rhs-modular}, and the negative sector follows from \eqref{eq:halfshift}.
\end{proof}

\subsection{The third through sixth companions}
The third through sixth companions admit certificates of the same kind as \eqref{eq:k2cert}. Because of their size, their coefficients are recorded in the accompanying data (\cref{app:verification}).

\begin{theorem}[Third through sixth companions]\label{thm:k3to6}
Let $k\in\{3,4,5,6\}$ and $m\ge k$. Identity \eqref{eq:wron} holds for both $z=0$ and $z=\frac12$. Consequently, the three pairs $(B(k,z),C(k,z))$ with $z=-\frac12,0,\frac12$ lie in $\M_{2m+1}$.
\end{theorem}

Retain $u=q^v$, $B_v=B(-2v\one,-v)$ and $E_i=e_{2i}-e_{2i-1}$. Define
\[
G_{m,k}(v;q)=F_{2m+1,\,B_v+\sum_{i=m-k+1}^mE_i,\,0}(q),
\]
and abbreviate $R_k(e)=R_k(e;w/u)$ from \eqref{eq:Rdef}.

\begin{lemma}[Certified boundary reductions]\label{lem:certred}
Suppose either $k\in\{3,4\}$ and $v\in\Q$, or $k\in\{5,6\}$ and $v\in\{0,\frac12\}$. For every $m\ge k$,
\begin{equation}\label{eq:certred}
G_{m,k}(v;q)=q^{-kv-k^2/2}(-q^{1/2-v};q)_\infty\sum_{j=0}^kA_{k,j}(q)\,\HL_m(q^j;u\one;q).
\end{equation}
\end{lemma}

\begin{proof}
Only the last $2k+1$ coordinates occur. Write $a_i=q^{n_{r-2k+i}}$ for $0\le i\le2k$, and put $M_k=\prod_{i=0}^{2k}a_i$. In the algebraic calculation the $a_i$, $w$ and $u$ are independent indeterminates when $k=3,4$; for $k=5,6$ we specialize $u=1$ or $u=w$. Let $\sigma_i$ replace $a_i$ by $qa_i$, with $q=w^2$, and define
\begin{align*}
\mathcal T_iP&=a_{i-1}a_{i+1}(1-a_i)P-u^{2(-1)^{i+1}}a_i^2\sigma_iP\qquad(1\le i<2k),\\
\mathcal T_{2k}P&=a_{2k-1}(1-a_{2k})P-(wu)^{-1}a_{2k}\sigma_{2k}P.
\end{align*}
For $k=3,4$, the files \file{uniform_k3_certificate.json} and \file{uniform_k4_certificate.json} define Laurent polynomials $P_1,\dots,P_{2k}$ as finite coefficient lists: an entry with index $i$, exponent vector $\alpha$ and coefficient $c(w,u)$ contributes $c(w,u)\prod_{h=0}^{2k}a_h^{\alpha_h}$ to $P_i$. For $k=5$ the corresponding lists are \file{uniform_k5_zero_certificate.json} at $u=1$ and \file{uniform_k5_half_certificate.json} at $u=w$; their coefficients are Laurent polynomials in $w$ alone. For $k=6$, the files \file{uniform_k6_zero_certificate.json} and \file{uniform_k6_half_certificate.json} contain integer polynomials in $w$, and each also specifies, in its field \file{target_multiplier}, a nonzero polynomial $D_{6,v}(w)$. Set $D_{k,v}=1$ for $k=3,4,5$. The coefficient lists satisfy the cleared identity
\begin{equation}\label{eq:certid}
\sum_{i=1}^{2k}\mathcal T_iP_i=D_{k,v}(w)\Bigl(w^{k^2}u^ka_{2k}\prod_{h=0}^{k-1}a_{2h+1}^2-M_kR_k(a_{2k})\Bigr).
\end{equation}
Identity \eqref{eq:certid} is checked by exact expansion in the Laurent polynomial ring over $\Z$, with all remaining variables independent; for $k=6$ the expansion is carried out one tail monomial at a time, and $D_{6,v}\ne0$ is checked directly.

Define $\Lambda$ as in the proof of \cref{lem:k2red}, using these $2k+1$ tail coordinates and the base shift $B_v$. After conjugation by $M_k^{-1}$, the reindexing factors are $u^{2(-1)^{i+1}}$ at an interior coordinate and $(wu)^{-1}$ at the last coordinate. Hence $\Lambda(M_k^{-1}\mathcal T_iP_i)=0$. The factors $1-a_i$ remove the lower boundary terms, and positive definiteness permits the shifts in the absolutely convergent sums. Applying $\Lambda(M_k^{-1}\,\cdot\,)$ to \eqref{eq:certid} and cancelling $D_{k,v}(w)$ wherever it is nonzero gives
\[
q^{kv+k^2/2}G_{m,k}(v;q)=\sum_{j=0}^kA_{k,j}(q)\,(-q^{1/2-v};q)_j\,F_{r,B_v+je_r,0}(q).
\]
Both sides are holomorphic on the upper half-plane, so the identity theorem extends the equality across the zeros of $D_{k,v}(w)$. By \cref{lem:boundary},
\[
F_{r,B_v+je_r,0}(q)=(-q^{j+1/2-v};q)_\infty\HL_m(q^j;u\one;q).
\]
Combining the finite and infinite products gives \eqref{eq:certred}.
\end{proof}

\begin{proof}[Proof of \cref{thm:k3to6}]
Combining \cref{lem:certred,lem:inverse} yields
\[
G_{m,k}(v;q)=q^{-kv}(-q^{1/2-v};q)_\infty J_k(q^v\one).
\]
At $v=0$, the confluent limit \eqref{eq:conf0} gives \eqref{eq:wron} for $z=0$. At $v=\frac12$, the half-shift identity \eqref{eq:halfshift} and $(-1;q)_\infty=2(-q;q)_\infty$ give
\[
S^{(1/2)}_{m,k}=q^{k/2}(-q;q)_\infty J_k(w\one).
\]
The factor $q^{-k/2}$ in \eqref{eq:confhalf} cancels the displayed power of $q$. \Cref{prop:rhs-modular} proves modularity, and \eqref{eq:halfshift} supplies the negative half-shift sector.
\end{proof}
